\documentclass[sigconf]{acmart}

\AtBeginDocument{%
  \providecommand\BibTeX{{%
    Bib\TeX}}}

\definecolor{paperblue}{HTML}{7f99b0}
\newcommand{\guideline}[1]{\textit{\color{paperblue}{#1}}}


\paperID{UNCIRCUIT} % your project ID from your proposal

\begin{document}


\title{Exploring the Evidence and Limits of Mechanistic Explanations}




\author{Mauricio Villavicencio}
\email{mivilla@umd.edu}
\affiliation{%
 \institution{University of Maryland, College Park}
  \city{}
  \state{}
  \country{}
}

\author{Bao Ngo}
\email{baongo85@umd.edu}
\affiliation{%
 \institution{University of Maryland, College Park}
  \city{}
  \state{}
  \country{}
}

\author{Parthav Poudel}
\email{baongo85@umd.edu}
\affiliation{%
 \institution{University of Maryland, College Park}
  \city{}
  \state{}
  \country{}
}

\author{Belu Ticona}
\email{ticona.belu@gmail.com}
\affiliation{%
 \institution{George Mason University}
  \city{}
  \state{}
  \country{}
}



\begin{abstract}

\guideline{DRAFT abstract, 210 words. Rewrite the last two sentences once the results section is final.}

\noindent Language model interfaces report confidence from the output distribution. That number is the
probability of the answer the model produced, and it says nothing about how stable the mechanism behind the
answer is. A recent line of work measures that stability instead, by asking how large an internal intervention
is needed to change the answer. We ask how such a fragility signal should be shown to people. We built a
prototype over 24 short factual questions answered by Qwen2.5-0.5B-Instruct and computed an approximate causal
robustness margin for each answer. The same signal is rendered five ways: baseline output confidence, a single
banded margin score, a per-layer flip curve, an interactive push-it-yourself slider, and a plain-language
counterfactual sentence. An overview scatter plots confidence against margin across the question set.
Seventeen participants worked through a self-paced session with graded reading tasks, a prediction task,
per-view ratings and a comprehension check. We report how accurately people read each representation and which
views they ranked as most useful. We also report whether they could tell a precomputed slider from a live
model, and where the vocabulary of the design broke down against the numbers on screen. The study is small and
exploratory. We read it as design evidence rather than as a test of effects.

\noindent \textbf{Sharing preference}: This paper may be shared [publicly | with future 839E students only (default) | not at all].
\end{abstract}




\keywords{Mechanistic interpretability, uncertainty communication, interactive explanation, activation
steering, visualization design, user study}



\begin{teaserfigure}
  \includegraphics[width=\textwidth]{img/teaser.png}
  \caption{Teaser example, although you don't have to use a teaser.}
  \label{fig:teaser}
\end{teaserfigure}



\maketitle



\section{Introduction}

\guideline{Outline to write next. Target 1.5 columns.}
\begin{itemize}
  \item Products report output confidence. That number says nothing about whether the mechanism behind the
    answer is stable. Name the gap in one paragraph.
  \item Mechanistic uncertainty: the causal robustness margin of
    \citet{komanduriMeasuringMechanisticUncertainty2026}, the smallest internal intervention that changes the
    output. Incorrect and uncertain answers have smaller margins.
  \item The question of this project is not whether the signal can be measured, but how it should be shown.
    State the three proposal questions, narrowed to the representation one.
  \item Contributions. One prototype that renders a single mechanistic signal five ways over the same question
    set. An exploratory study with 17 participants on reading accuracy, perceived usefulness and comprehension.
    A set of design problems the study surfaced, chiefly the scale-sense problem and the confusion between a
    live model and a precomputed replay.
  \item Say plainly here that our margin approximates the published method. On our item set it is strongly
    rank-correlated with output confidence, $\rho = 0.71$. That is a limit of the stimulus, and the paper is
    about representation either way.
\end{itemize}


\section{Related Work}

Three bodies of work meet in this project. Interactive interfaces intervene on a model's internals. Visual
analytics shows how fragile a prediction is. Uncertainty communication studies how such a signal should be
drawn or worded for a reader. As far as we can tell they have not met before. We take each in turn and then
state where our work sits.

\subsection{Interfaces that intervene on model internals}

Visual analytics for neural models moved early from passive inspection to intervention. LSTMVis let analysts
select a span of text as a hypothesis about hidden state and match it against a corpus
\cite{strobeltLSTMVisToolVisual2016}. Seq2Seq-Vis then added what has become the canonical interaction for the
genre: change the model's internal state, and watch the output change
\cite{strobeltSeq2SeqVisVisualDebugging2018}.

The same move appears in model-agnostic tools. Prospector let analysts drag a feature value on a single
instance and read the prediction off a partial-dependence curve \cite{krauseInteractingPredictionsVisual2016}.
The What-If Tool made edit-a-datapoint-and-re-infer a standard idiom \cite{wexlerWhatIfToolInteractive2019}.
HypperSteer applied interactive perturbation to sequence models for hypothetical testing
\cite{wangHypperSteerHypotheticalSteering2020}.

For transformer language models, LM-Debugger is the closest existing system to our interactive view. It
interprets feed-forward updates in vocabulary space and lets a user configure interventions that shift the
prediction in a chosen direction \cite{gevaLMDebuggerInteractiveTool2022}. The LM Transparency Tool exposes the
input-to-output information flow with attribution to individual heads and neurons
\cite{tufanovLMTransparencyTool2024}. AttentionViz offers the global-view-plus-detail structure we borrow for
our overview-to-item drill-down \cite{yehAttentionVizGlobalView2023}. All of these are built for experts, and
they ask what a component does rather than how sure the model is. Two exceptions matter to us. Transformer
Explainer explains internals to non-experts through live interaction
\cite{choTransformerExplainerLearning2024}. Steerable Chatbots hands an end user a steering coefficient as a
slider and studies how much agency people want over it \cite{boSteerableChatbotsExploring2025}.

Why interaction at all? \citet{chengExplainingDecisionMakingAlgorithms2019} ran the comparison directly. In an
algorithmic-admissions setting, interactive explanations improved both objective and self-reported
understanding, while static white-box explanations improved only objective understanding. Interactivity also
cost participants time. \citet{rossEvaluatingInterpretabilityGenerative2021a} go further and use interactive
manipulation as the measurement of interpretability rather than as the interface under test, which shaped our
task design. \citet{kirchlerExplainabilityRequiresInteractivity2021} argue from the model side that static
attributions under-determine behaviour, so counterfactual interaction is required. Our push-it-yourself view is
that tradition applied to a fragility signal. The slider does not tune a feature or a personality. It spends
the margin.

\subsection{Showing fragility as a response to perturbation strength}

A second literature draws fragility as a curve over perturbation magnitude, almost always in input space.
Bluff visualizes activation pathways over a range of attack strengths
\cite{dasBluffInteractivelyDeciphering2020}. AdvEx gives novices a perturbation adjuster and reports accuracy
and confidence before and after the attack, which makes it our nearest neighbour in both encoding and audience
\cite{youPandaNotPanda2023}. \citet{liVisualAnalyticsNeuron2023} ask which units are most vulnerable, and that
is the per-layer question our layer-by-strength inset answers. On the modelling side,
\citet{stoehrActivationScalingSteering2024} treat the scaling factor needed to flip a prediction as the
quantity of interest. It is the computational sibling of the margin, and the reason a strength sweep is a
defensible thing to plot. What none of this work does is treat the curve as an uncertainty display for a lay
reader. It remains a debugging and robustness-analysis idiom.

\subsection{Representing uncertainty for a lay reader}

Our single-score view inherits a long argument about how to encode uncertainty honestly.
\citet{schulzDeepViewVisualizingClassification2019} visualize classification boundaries together with per-point
certainty, which is the spatial version of how far an answer is from flipping.
\citet{padillaUncertaintyVisualization2021} survey the design space our manipulation samples from.
\citet{kayWhenIshMy2016} show that lay readers reason better about quantified uncertainty in discrete,
frequency-framed displays than in abstract intervals. \citet{fernandesUncertaintyDisplaysUsing2018} demonstrate
the lesson we took most seriously. The display that best supports estimating a quantity is not necessarily the
one that best supports deciding with it. That is why we measure prediction accuracy and not only comprehension.
\citet{kaleVisualReasoningStrategies2020} repeat the result with a strategy-level qualitative analysis, which is
a model for how we coded open responses.
\citet{correllValueSuppressingUncertaintyPalettes2018} address the specific risk of our margin score, a derived
number presented as if it were precise. \citet{bisantzDisplayingUncertaintyInvestigating2005} report the
preference-performance split we see in our own ratings, where the familiar bar is comprehended well but is not
judged most helpful.

\subsection{Counterfactual explanations in plain language}

Our sentence view is a counterfactual explanation with the intervention in internal rather than input space.
DECE treats counterfactuals as the user-facing unit and adds interactions for customizing them
\cite{chengDECEDecisionExplorer2020}. GAM Coach has users slide feature values under their own constraints until
the prediction flips, which is our interactive view and our sentence view fused into one interaction
\cite{wangGAMCoachInteractive2023}. \citet{schlegelInteractiveCounterfactualGeneration2024} extend interactive
counterfactual authoring beyond tabular data. On whether the format helps,
\citet{kuhlLetsGoAlien2023} provide a controlled framework and evidence that counterfactual explanations beat no
explanation on objective performance and subjective usability.
\citet{kuhlKeepYourFriends2022} show that the closest counterfactual teaches better than a merely plausible one.
That is why our sentence reports the minimum push rather than a vivid one.

\subsection{Communicating LLM uncertainty, and reliance}

Our baseline condition is the signal this literature studies. \citet{kimImNotSure2024} found, with 404
participants, that first-person expressions of uncertainty lowered agreement with the system and raised
accuracy. \citet{boRelyNotRely2024} temper what such interventions achieve. They reduce over-reliance but
largely fail to produce appropriate reliance. That sets the bar for any new display, and it warns us against
reading a preference for one of our views as a trust improvement.
\citet{villavicencioNotAllUncertainty2026} show that the granularity at which uncertainty is expressed changes
verification behaviour and self-confidence. \citet{warrenWhenShouldTrust2026} compare indicator formats
directly. All of these measure the output distribution or its verbalization. The internal mechanism is not part
of the signal.

Two further strands informed the study rather than the views. Behavioural evaluation tools frame model weakness
as something to slice and prioritize \cite{cabreraZenoInteractiveFramework2023,robertsonAnglerHelpingMachine2023},
which is how our overview quadrants are meant to be read.
\citet{cappuccioExplanationUserInterfaces2025} survey explanation user interfaces as a design space, which we
used to check that our five views do not miss a family.

\subsection{Where this work sits}

The premise that internal instability is worth reporting is not ours.
\citet{marjanovicInvestigatingImpactModel2024} show that explanations shift under perturbations that leave
predictive uncertainty looking confident. What is missing is the representation question. Interactive
interpretability tools intervene on internals but never report fragility to a reader. Fragility visual analytics
reports it in input space, for practitioners. Uncertainty communication studies the reader, but only ever with
the output distribution in hand. We treat a mechanistic fragility signal as an uncertainty display, render it
five ways over one question set, and put the comparison in front of non-expert participants.


\section{Methods}

\subsection{Stimulus: an approximate robustness margin}

We follow the causal robustness margin of \citet{komanduriMeasuringMechanisticUncertainty2026}, the smallest
internal intervention that changes a model's output. Our implementation approximates it, and the prototype says
so in an ``About this method'' panel.

We used \texttt{Qwen2.5-0.5B-Instruct}, which has 24 decoder layers, over 24 short factual questions with
single-token answers. The items span easy, medium and hard difficulty. The model answers in a few-shot,
one-word format, and we read logits at the last prompt position. For each item we record the answer token $a$
and its probability $p_{\mathrm{answer}}$, which is the baseline output confidence. We also record the
highest-probability alternative token $b$ whose decoding differs from $a$.

The intervention is a contrastive logit-lens direction, $v = W_U[b] - W_U[a]$, normalized to $\hat{v}$. This is
the special case of the Jacobian lens used in the original work with $J = I$. At decoder layer $\ell$, and at
the last position only, we add $\alpha \cdot \lVert h_\ell \rVert \cdot \hat{v}$ to the layer's output hidden
state. Scaling by the residual norm makes $\alpha$ unitless and comparable across layers. Sweeping 40 values of
$\alpha$ up to $1.5$ for every layer gives, per layer, the probability of $a$, the probability of $b$, and the
current top-1 token.

The per-layer margin is the smallest $\alpha$ at which top-1 is no longer $a$, refined by binary search between
bracketing grid points. Tokens that are merely a different surface form of the same answer do not count as a
flip. The margin for an item is the minimum over layers. We also record the layer that achieves it and the
item's percentile within the set.

Two properties of this stimulus matter for reading our results. First, every answer flips well below the top of
the sweep, and margins are small in absolute terms, from 0.003 to 0.077. This is the source of the scale-sense
problem we report. Second, the margin is strongly rank-correlated with output confidence on this set, with
Spearman $\rho = 0.71$, so the two signals do not diverge as sharply as the design assumes. The prototype
reports $\rho$ on the overview screen rather than hiding it. Twenty-one of the 24 answers are correct.

\subsection{The five representations}

The same precomputed data is rendered five ways. A fixed colour marks the answer token throughout, and a
contrasting colour marks the alternative.

\begin{description}
  \item[R0 Baseline confidence] the output probability of the answer as a bar, with the top-5 token
    distribution. This is what products show today.
  \item[R1 Margin score] the margin on a gauge with three bands, Fragile, Moderate and Robust, set at tertiles
    of the item set. It carries a numeric label and one plain-language comparison line.
  \item[R2 Flip curve] probability of the answer against push strength, one line per layer. The margin layer is
    highlighted, a marker sits at the flip, and an inset shows layer by strength, coloured by which token is
    top-1.
  \item[R3 Push it yourself] two sliders, strength and layer, over the precomputed grid. It shows a live readout
    of what the model would answer, animated probability bars, and an autoplay sweep.
  \item[R4 Counterfactual sentence] one sentence naming the size of the nudge, the target token, the strength
    and the layer, plus a percentile comparison.
\end{description}

An overview scatter plots output confidence against margin for all 24 items, with median splits and quadrant
labels. That is where ``confident but fragile'' is meant to become visible.

\subsection{Design}

The study is within-subjects, so every participant sees all five representations. The views are also the
material participants are asked to rate. Exposure order is therefore counterbalanced per participant for
R1--R4 in the guided part, reused for the rating section, and recorded with each response as
\texttt{view\_order}. Items for the graded tasks come in two sets, A and B, assigned at random, so that
participants working next to each other do not converge on the same answers. Set assignment is recorded as
\texttt{task\_set}.

\subsection{Pilot}

We ran a pilot round first, with about ten people, mostly classmates. Eight of them submitted complete
responses through the instrument. All eight were graduate students, five at master's level and three at PhD
level. The pilot
tested the instrument rather than measuring anything, and it changed the interface and the questionnaire enough
that its responses are not comparable with the main study's.

The pilot produced the following changes. A welcome page now states the goal, the duration, and that the session
can be paused. An outline of the five parts was added, so that participants are never dropped into the explorer
without framing. Free exploration became its own timed part, with the views hidden until the participant
starts. The former pop-up quick guide was folded into two framing steps at the start of the guided part.
Exposure to R1--R4 became counterbalanced, with the order recorded. The overview gained a linear axis with zoom
and a mini-map, after pilot participants misread the original scale. Each variable now has one colour across
the views, and the explanation text is plainer and larger. The answer key now says ``the model was
right/wrong'' instead of ``correct/incorrect''. The last build also lengthened the view-by-view walkthrough to
six to eight minutes.

For that reason the analysis below uses only sessions run on the final build. Every response records the build
it came from, so the split is auditable rather than a judgement call.

\subsection{Participants}

For the main study we recruited by convenience sampling. Participants came from an undergraduate NLP class and
from the authors' own contacts, who were free to pass the link on, so the sample also grew by snowball. Sessions
were remote and self-paced. Participation was voluntary, anonymous and uncompensated. Every participant gave
consent on the first screen before seeing the interface, and no identifying information was collected.

Of 25 recorded sessions in total, we analyze the 17 collected on the final build, \texttt{v0.5.0}, between 19
and 23 November 2026. The other 8 are the pilot responses described above. Of the 17, 7 were assigned task set
A and 10 set B. Median completion time was 31 minutes against a 20-minute design target.

We collected no age and no gender. The background items ask for program level and for four kinds of relevant
experience. The sample is mostly undergraduate, with 15 of 17 participants (88\%) in an undergraduate program,
one master's student and one who selected ``other''. Self-rated familiarity with large language models, on a
scale of 1 to 5, centres on the midpoint. The median is 3, with 11 of 17 (65\%) choosing exactly 3 and 5 (29\%)
choosing 4 or 5. These are people who use language models but do not build them.

Familiarity with mechanistic interpretability is low, which follows from the recruitment pool. The median is 1
and the mean 1.71, with 14 of 17 (82\%) at 1 or 2 and nobody at the top of the scale. Familiarity with
confidence or uncertainty scores in machine learning sits in between, at a median of 3 and a mean of 2.65, with
7 of 17 (41\%) at 1 or 2. We also asked whether participants had ever used a tool that edits or steers a
model's internal activations. Four said yes, 3 said no, and 10 (59\%) said they were not sure. More
participants were unsure than said no, which is itself a reading of how unfamiliar the concept is.

Two limits follow. The sample is self-selected and drawn largely from one class, so it is not representative.
We also did not collect any measure of visualization literacy or HCI background, which would have been the
natural covariate for how people read the five views. The four experience items are the only handle we have, so
we treat the study as exploratory throughout.

\subsection{Procedure}

The session runs in the browser in five parts. Each part is gated, so participants cannot read ahead.

\begin{enumerate}
  \item \textbf{Consent and background (A).} Consent, program level, and four familiarity items.
  \item \textbf{Free exploration, then first impressions (B).} The views stay hidden until the participant
    presses ``Start exploring''. The controls are highlighted, and clicking the overview selects an item rather
    than jumping to a view. First impressions unlock after three minutes. They ask what the interface shows,
    what ``Robustness'' means in the participant's own words, which views were opened, whether the About panel
    was read, and what first confused them. These answers capture the unprompted mental model, before any task
    shapes it.
  \item \textbf{Guided tasks (C).} Two framing steps come first. ``Learn the ideas'' is identical for everyone
    and folds in the former quick guide as ``Try a push'' and ``Two scores''. ``Explore the views'' then walks
    R1--R4 in the participant's counterbalanced order, over about six to eight minutes. The graded tasks cover
    reading one answer, which means its confidence, the token it flips to, and the easiest layer. They then
    cover steering with the slider and comparing the two signals. A prediction task follows, in which
    participants judge for four items whether the model's answer is right. Predictions are submitted before the
    answer key unlocks. The key then reports whether the model was right or wrong, and whether each prediction
    matched.
  \item \textbf{Steering experience (D) and per-view ratings (E).} Seven agreement items ask about the push
    interaction. Two more ask whether it felt like changing a live model or replaying recorded results, and
    whether the ``values are precomputed'' caption was noticed. Section E rates each of the five views on three
    statements and ranks them from most to least useful. It also probes the vocabulary of the design directly.
    One item asks whether a margin of 0.048 labelled ``Robust'' read as a small or a large number. Another asks
    whether calling that strength ``a large internal nudge'' matched what was on screen.
  \item \textbf{Understanding check (F), trust and open feedback (G, H).} Seven graded comprehension items ask
    what the bands are relative to, what the grey lines in the flip curve are, where the push is applied, and
    whether a surface-form change counts as a flip. Sections G and H collect agreement items on trust and
    usefulness, and free-text feedback.
\end{enumerate}

Section B receives no hints beyond the interface itself, so that the About panel remains part of what is being
tested. Sections F and D2 therefore measure comprehension after the guided framing. Participants never see
grades.

One planned task was dropped before data collection. No item in the final set combines high output confidence
with a Fragile-band margin, so the task that would have tested the ``confident but fragile'' quadrant directly
has no stimulus. The final comprehension item instead tests fragility under the intervention we actually ran,
rather than claiming that fragility predicts correctness.

\subsection{Measures}

\begin{itemize}
  \item \textbf{Reading accuracy.} Graded short answers in C and seven graded items in F. Each offers a
    ``Not sure'' option, so that guessing appears as uncertainty rather than as error.
  \item \textbf{Prediction.} For four items, a five-point judgement of whether the model's answer is right,
    recorded per item together with the views the participant used to decide.
  \item \textbf{Perceived usefulness and clarity.} Three 7-point agreement statements per view, and a forced
    ranking of the five views.
  \item \textbf{Calibration of the vocabulary.} The scale-sense items, E5 and E6, and the band question, F3.
  \item \textbf{Attribution of control.} Whether the slider felt live or precomputed (D2), and whether the
    precomputed caption was read (D3).
  \item \textbf{Trust and usability.} The 7-point agreement items in G, plus open responses in E7, D4, D5 and H.
  \item \textbf{Behaviour.} Interaction events logged by the instrument, including view switches, item
    selection, slider moves, autoplay use, About-panel opening, guide exposure, and per-step timings.
\end{itemize}

\subsection{Apparatus and instrumentation}

The prototype is a static web application. The margins are computed offline by a Python pipeline and shipped as
JSON, so nothing runs a model in the browser. This makes the interface instant and impossible to crash during a
session. The cost is a slider that replays precomputed values, which is what D2 and D3 probe.

The study instrument lives at its own address and saves progress locally, so a session can be resumed. It posts
one row per participant and one row per interaction event to a spreadsheet, and responses can also be downloaded
as CSV and JSON. Every response carries \texttt{app\_version}, \texttt{task\_set}, \texttt{view\_order} and
\texttt{guidance\_version}. Older builds stay online at their own addresses, so a version can be inspected after
the fact.

\subsection{Analysis}

We have 17 participants, and a within-subjects design whose exposure order we counterbalanced but did not vary
as a factor. We therefore report descriptive statistics and treat the study as exploratory. For every graded
question we report counts of correct, wrong and ``not sure''. For the rating statements we report per-view means
and distributions, and we report the ranking of the five views. Prediction accuracy is reported per item,
alongside that item's true band and correctness, and timing is reported as medians with quantiles. Scripts in
the repository compute these from the exported responses, so the numbers in the next section are reproducible
from the raw CSV.

Open responses, which are B1, B2, B6, D4, E7 and H, are coded by one author into themes, and a second author
reviews the assignments. We report themes with counts and representative quotations, following the
strategy-coding approach of \citet{kaleVisualReasoningStrategies2020}. We make no inferential claims about
differences between views.

\subsection{Ethics}

The consent screen stated that the study is part of a course project and takes about twenty minutes. It also
stated that participation is voluntary, and that responses are anonymous and used only for this report.
Participants could stop at any time, and declining on the first item ended the session. Responses carry no
identifying information. Use is limited to this report, and we would seek IRB review before publishing beyond
the course. Because some participants were recruited from a class, we took care that participation was not tied
to coursework or credit.


\section{Preliminary Results}

\guideline{Outline. All numbers below are already computed in \texttt{results/final\_study.json}. Write prose
against them, and add the two or three figures that carry the argument.}
\begin{itemize}
  \item Reading accuracy by task (C1--C10) and comprehension (F1--F7). Where accuracy is high, and the three
    items where it collapses. Report ``not sure'' separately from wrong.
  \item Prediction task. Per-item judgements against each item's band and true correctness. Note the one item
    where nearly everyone was right, and the one where the views did not help.
  \item Per-view ratings and the forced ranking. Which views won, and whether the interactive view's advantage
    survives the ranking.
  \item Scale sense. How many read 0.048 as small, and whether ``a large internal nudge'' matched.
  \item Live-versus-precomputed attribution, against whether the caption was noticed.
  \item Timing and interaction logs. Median 31 minutes against a 20-minute target, and which views were opened
    unprompted in B.
  \item First-impression themes before any framing. What people thought ``Robustness'' meant.
\end{itemize}


\section{Discussion}

\guideline{Outline.}
\begin{itemize}
  \item The representation question is separable from the measurement question. Our results speak only to the
    first.
  \item Scale sense is a design problem for any derived fragility number. Bands relative to a question set make
    a tiny absolute number read as ``Robust'', and a small push read as ``large''.
  \item Interaction buys comprehension, but it also invites over-attribution of control. Relate to
    \citet{chengExplainingDecisionMakingAlgorithms2019} and \citet{boSteerableChatbotsExploring2025}.
  \item What $\rho = 0.71$ means for the pitch. On this item set the mechanistic signal is largely redundant
    with output confidence, so the claim needs an item set where the two diverge.
  \item Reliance framing. Following \citet{boRelyNotRely2024}, a preference for a view is not evidence of
    appropriate reliance.
\end{itemize}


\section{Future Work}

\guideline{Because this is a course project, substantial future work is acceptable.}

\begin{itemize}
  \item An item set selected for divergence between output confidence and margin. That is the stimulus this
    design needs, and the one we did not have.
  \item The paper's Jacobian-lens direction and a larger model, instead of the contrastive logit-lens
    approximation.
  \item Absolute rather than set-relative bands, and a vocabulary test for the words ``fragile'' and ``nudge''.
  \item A live-computation slider, to separate the interaction from the precomputed-replay artifact.
  \item A powered between-subjects comparison of two or three views on a decision task, with appropriate
    reliance as the dependent measure. It should draw on a sample beyond one class and collect a
    visualization-literacy measure as a covariate.
\end{itemize}


\section*{Author contributions}

\guideline{
You can refer to \href{http://li.mit.edu/Archive/Activities/PubFormat/Nature/End_notes.pdf}{\underline{Nature's annotated example}}.
Our example: Testudo Turtle wrote all the code. Seaweed Grass debugged the code and analyzed all the data. Both
authors designed the experiment, discussed the results and future work, and prepared the manuscript.}


\section*{Acknowledgment}

\guideline{Optional}

\section*{Supplementary materials}

\guideline{Recommended. Candidates: the live prototype and the study address, with older versions kept online,
one per version. The questionnaire source. The offline pipeline and its README, which document the
approximation and $\rho$. The response CSV and the analysis script. The annotated related-work map used to
position this paper.}


\bibliographystyle{ACM-Reference-Format}
\bibliography{refs,sample-base}

\end{document}
